% ===============================
% SEZIONE 1: Esercizi
% ===============================
\begin{question}[ID=1.1, subtitle={Completamento con percentuali — serie 1}, tags={percentuali}, difficulty=facile]
Completa con i numeri opportuni:
\begin{center}\begin{minipage}{0.92\linewidth}
\begin{multicols}{2}
\begin{enumerate}
\item $\text{Il }20\%\text{ di }150\text{ è } \rule{1.5cm}{0.4pt}$
\item $\text{Il }15\%\text{ di }200\text{ è } \rule{1.5cm}{0.4pt}$
\item $\text{Il }18\%\text{ di } \rule{1.5cm}{0.4pt} \text{ è }3{,}6$
\item $\text{Il } \rule{0.5cm}{0.4pt}\%\text{ di }80\text{ è }4$
\end{enumerate}
\end{multicols}
\end{minipage}\end{center}
\begin{solution}$30$; $30$; $20$; $5\%$\end{solution}
\end{question}

%--------------------------------------------------%

\begin{question}[ID=1.2, subtitle={Completamento con percentuali — serie 2}, tags={percentuali}, difficulty=facile]
Completa con i numeri opportuni:
\begin{center}\begin{minipage}{0.92\linewidth}
\begin{multicols}{2}
\begin{enumerate}
\item $\text{Il }25\%\text{ di }160\text{ è } \rule{1.5cm}{0.4pt}$
\item $\text{Il }30\%\text{ di }90\text{ è } \rule{1.5cm}{0.4pt}$
\item $\text{Il }12\%\text{ di } \rule{1.5cm}{0.4pt} \text{ è }6$
\item $\text{Il } \rule{0.5cm}{0.4pt}\%\text{ di }250\text{ è }50$
\end{enumerate}
\end{multicols}
\end{minipage}\end{center}
\begin{solution}$40$; $27$; $50$; $20\%$\end{solution}
\end{question}



% ===============================
% SEZIONE 2: Problemi
% ===============================
\begin{question}[ID=2.1, subtitle={Prezzo scontato del maglione natalizio}, tags={percentuali}, difficulty=facile]
Quest'anno per Natale ho deciso: solo regali utili. Un negozio di abbigliamento applica uno sconto del $\text{25\%}$ su un maglione a tema natalizio.

\begin{enumerate}
    \item Sapendo che il prezzo iniziale del maglione è di $\text{39,99}$ euro, determina il prezzo scontato.
    \color{gray}\item In un altro negozio, lo stesso maglione viene venduto a $\text{25}$ euro già scontato del $\text{20 \%}$. Determina il prezzo iniziale del maglione prima dell'applicazione dello sconto.
\end{enumerate}
\begin{solution}1) $39{,}99\times0{,}75\approx 29{,}99$ €; 2) $25\div0{,}80=31{,}25$ €\end{solution}
\end{question}


%====================================================%
% SEZIONE 3: Notazione scientifica
%====================================================%
\begin{question}[ID=3.1, subtitle={Conversione in notazione scientifica}, tags={notazione scientifica}, difficulty=facile]
Trasforma le seguenti misure in notazione scientifica:
\begin{center}\begin{minipage}{0.92\linewidth}
\begin{multicols}{2}
\begin{enumerate}
\item $0{,}0045 = \rule{3cm}{0.4pt}$
\item $32000 = \rule{3cm}{0.4pt}$
\item $0{,}00000082 = \rule{3cm}{0.4pt}$
\item $5\,600\,000 = \rule{3cm}{0.4pt}$
\item $123{,}45 = \rule{3cm}{0.4pt}$
\item $34{,}5 \cdot 10^{-2} = \rule{3cm}{0.4pt}$
\end{enumerate}
\end{multicols}
\end{minipage}\end{center}
\begin{solution}$4{,}5\times10^{-3}$; $3{,}2\times10^4$; $8{,}2\times10^{-7}$; $5{,}6\times10^6$; $1{,}2345\times10^2$; $3{,}45\times10^{-1}$\end{solution}
\end{question}
